% =============================================================
% ProcessHeal — Paper (revised: adds Goals/Objectives + the
%   commercial-BAS-FDD distinction; tightened intro)
% UREAP 2026, Thompson Rivers University
% Author: Yassh Singh
% Supervisors: Dr. Anthony Aighobahi, Dr. Mridula Sharma
%
% Version history (all in this folder):
%   paper-v1.tex      — first pass, pre-critique
%   paper-draft.tex   — earlier paper-size version
%   paper.tex         — THIS FILE, current working draft
%   literature-review.tex — long survey for the UREAP final report
%
% Compile in Overleaf: Recompile twice. Bibliography in
% references.bib (shared across all .tex files in this folder).
% =============================================================

\documentclass[11pt,letterpaper]{article}

\usepackage[margin=1in]{geometry}
\usepackage{microtype}
\usepackage[hidelinks]{hyperref}
\usepackage{cite}
\usepackage{times}

\title{ProcessHeal: Process Mining for Fault Detection and Diagnosis in Building HVAC Systems}
\author{Yassh Singh \\
\small Supervisors: Dr.\ Anthony Aighobahi, Dr.\ Mridula Sharma \\
\small Thompson Rivers University --- UREAP 2026}
\date{May 2026}

\begin{document}
\maketitle

\begin{abstract}
\noindent
Heating, ventilation, and air conditioning (HVAC) systems account for a substantial share of building energy use, and 15--30\% of that energy is wasted by faults that operate quietly below the alarm threshold of a building's own control system. Recent empirical work has shown that seven of the ten most common air-handling-unit faults are behaviour-based rather than condition-based, which suggests that detection methods built around modelling temporal ordering have a structural advantage over the static classifiers and rule libraries that dominate current practice. We propose \emph{ProcessHeal}, a pipeline that applies process mining to building HVAC fault detection. The pipeline learns a Petri-net model of expected control behaviour from a building's own healthy operating data and uses alignment-based conformance checking to identify and localise deviations to specific equipment. The discovered Petri net is the diagnostic explanation, in contrast to machine-learning approaches that pair black-box classifiers with post-hoc attribution. Because the expected behaviour is learned rather than pre-programmed, ProcessHeal complements rather than duplicates the fixed rule libraries shipped with commercial building automation systems, which detect only the fault patterns coded into them in advance. To our knowledge, this is the first application of event-log process mining to building HVAC, a claim corroborated by four independent systematic reviews. We validate the pipeline on the publicly available LBNL Single-Duct AHU benchmark and outline a transfer to Thompson Rivers University's Campus Activity Centre, producing a recommissioning diagnostic report for facility operators as the deployment artefact.
\end{abstract}

\section{Introduction}

Buildings account for approximately 34\% of global energy-related carbon dioxide emissions, and HVAC systems carry the dominant share within that total \cite{bi2024aihvacreview}. The fraction of HVAC energy wasted by undetected faults has consistently been estimated at 15--30\% across recent reviews \cite{bi2024aihvacreview,pritoni2022faultcorrection}, but until late 2023 these numbers rested on relatively small studies. The first large-scale empirical characterisation appeared when Crowe and colleagues analysed multi-year fault-detection output across more than 60{,}000 pieces of HVAC equipment in 317 commercial buildings \cite{crowe2023faultprevalence}. Their study reports an average of 245 faults per building per month and finds 40\% of all air handling units in a faulted state on any given day. The most common AHU fault in their dataset, a supply-air-temperature setpoint that is never reset, was observed on 55\% of the AHUs surveyed. These buildings were not unmonitored; many run modern building automation systems with alarms and fault-detection features, yet the faults persisted, which points to a gap in the kind of fault these tools can see rather than in their deployment.

The nature of that gap is specific. Of the ten most common AHU faults catalogued by Crowe and colleagues, seven are behaviour-based rather than condition-based: incorrect schedules, unreleased manual overrides, missed setpoint resets, and broken economiser sequences manifest as deviations in the temporal ordering of system states rather than as isolated component failures. A threshold alarm fires when a single variable leaves its acceptable band, and a pre-programmed rule fires when a named fault pattern is matched, but neither mechanism can easily detect a control sequence that executes its steps in the wrong order or skips a step entirely. The structural argument that follows is straightforward: detection methods that model the sequence in which a building moves between operating states have a natural advantage on the fault classes that dominate the empirical landscape.

The dominant research methodology for HVAC fault detection does not model that sequence either. Machine learning has been surveyed across 211 AI-based HVAC FDD studies published between 2013 and 2023 \cite{bi2024aihvacreview}, with reported accuracies on the standard benchmarks saturating at 97--99\% across hybrid random-forest pipelines on RP-1312 data \cite{tun2021hybridrfsvm} and convolutional-and-LSTM hybrids on real-building data \cite{prince2025improved1dcnnlstm}. The most recent work has shifted toward interpretability, but in a specific way: a deep model is paired with a post-hoc explanation method such as layer-wise relevance propagation or attention attribution \cite{deng2024afgcn,wang2025xmhac}. The resulting explanation is at the level of input features rather than at the level of the control sequences the building was supposed to follow. A facility manager told that sensor $X$ had high relevance for a prediction knows that something is unusual about that signal, but is not told how the building was supposed to be behaving instead.

A second problem concerns reproducibility. A recent audit of 65 primary machine-learning HVAC fault detection studies found that only two shared executable code, and that 72\% do not specify whether their dataset is public, proprietary, or commercial \cite{mukhtar2025reproducibility}. The field has plateaued on a metric that may not be the right metric, with little of the existing work independently verifiable, and with the dominant explanation style answering a question (which signal mattered) different from the one operators ask (what should the system have been doing).

The opportunity is to change the model class. Process mining is a methodology in which a discrete event log derived from system observations is used to discover a Petri-net model of expected behaviour, with conformance checking against that model producing an alignment-based deviation signal that carries formal fitness and precision semantics \cite{vanderaalst2016processmining}. We apply this approach to HVAC fault detection, treating sensor-derived events as the input log, ASHRAE Guideline 36 control sequences as the conceptual reference, and the discovered Petri net as both the expected-behaviour model and the diagnostic explanation. Three contributions follow. ProcessHeal is, to our knowledge, the first application of event-log process mining to building HVAC fault detection, a novelty claim corroborated by four independent systematic reviews. The pipeline is evaluated on the publicly available LBNL Single-Duct AHU benchmark, the same dataset used by the strongest recent interpretable-machine-learning work, allowing direct comparison on equivalent fault scenarios. A diagnostic report for facility operators is produced as the deployment artefact, drawing on a planned transfer of the pipeline to real-world data from Thompson Rivers University's Campus Activity Centre, where it is intended to complement the rule-based fault detection already present in the building's automation system.

\section{Research Aim and Objectives}
\label{sec:aims}

The aim of this research is to establish process mining as an interpretable and generalisable approach to fault detection and diagnosis in building HVAC systems. The approach is first validated on a public benchmark dataset, then demonstrated on a real building. Six objectives carry the work from method to deployment:

\begin{enumerate}
\item \textbf{Abstract events from sensor data.} Convert continuous HVAC sensor streams into discrete event logs, using control-sequence rules grounded in ASHRAE Guideline 36.
\item \textbf{Discover and check the model.} Learn a Petri-net model of expected control behaviour from fault-free data, and detect faults as conformance deviations localised to specific equipment.
\item \textbf{Benchmark against known answers.} Evaluate the pipeline on the LBNL Single-Duct AHU dataset and its documented ground-truth faults, comparing accuracy and interpretability with published machine-learning methods on the same data.
\item \textbf{Test a hybrid extension.} Determine whether combining process-mining features with machine-learning classifiers detects faults more reliably than either method alone.
\item \textbf{Deploy on a real building.} Apply the validated pipeline to live trend data from a Thompson Rivers University building, delivering a ranked diagnostic report that complements the rule-based fault detection already present in the building's automation system.
\item \textbf{Release the work openly.} Publish the pipeline and its building configuration as open-source software, evaluated on a public dataset, to address the reproducibility gap identified in the literature.
\end{enumerate}

\noindent Objectives one to four are pursued on the LBNL benchmark as a controlled testbed; objective five constitutes the real-world deployment phase; objective six runs throughout the project.

\section{Related Work}
\label{sec:related}

\subsection{Machine learning for HVAC fault detection, and the move from accuracy to interpretability}

The machine-learning era of HVAC fault detection is well documented. Bi and colleagues survey 211 AI-based HVAC FDD studies published between 2013 and 2023, organising them into traditional machine-learning, deep-learning, hybrid AI, and physical-model branches \cite{bi2024aihvacreview}. Reported accuracies on the standard benchmark datasets have effectively saturated. Tun and colleagues reached 98\% accuracy across thirteen fault types using a hybrid random-forest and support-vector-machine pipeline on RP-1312 AHU data \cite{tun2021hybridrfsvm}. Prince, Yoon, and Kumar reached 97\% accuracy on data from three real buildings using a thirteen-layer one-dimensional convolutional-and-LSTM hybrid that remains robust to Gaussian noise and to 10\% missing values \cite{prince2025improved1dcnnlstm}. The authors of that paper close by explicitly identifying interpretability as the problem they did not solve, proposing grafting SHAP or LIME as future work.

The most directly relevant interpretable benchmark for the LBNL SDAHU subset is Deng and colleagues' adaptive fusion graph convolutional network. The model reaches 88.09\% accuracy at 500 labelled samples and 75.15\% at only seven samples, with per-node sensor importance as the explanation mechanism \cite{deng2024afgcn}. Subsequent work in the same vein pairs multi-head attention with parallel dilated convolutions and applies layer-wise relevance propagation for explanation on chiller fault data, reaching an F1 score of 0.9841 across 27 operating conditions \cite{wang2025xmhac}. The pattern across this body of work is consistent: a deep black-box network is the model, a post-hoc analysis layer provides the explanation, and the explanation is framed at the level of input features rather than control flow. The architectural opportunity is for a method whose model is its explanation.

\subsection{Unsupervised methods and cross-building deployment}

Real buildings rarely come with ground-truth fault labels, and the strongest recent work on unlabelled data uses self-supervised or autoencoder methods. Abdollah and colleagues pretrained a Transformer encoder via masked-token reconstruction on unlabelled HVAC time-series from a heat-pump building at the Polytechnic University of Milan, then used Peak-Over-Threshold detection on reconstruction error to flag a real two-month monitoring fault between October and December \cite{abdollah2024transformerssl}. The most directly relevant precedent for a Canadian-university real-BAS deployment is El Mokhtari and McArthur's autoencoder-based work at Toronto Metropolitan University, which distinguishes equipment-level from system-level faults across five fan coil units and demonstrates cross-unit transfer learning \cite{elmokhtari2024aebas}. These methods share a common limitation. Each produces a learned latent embedding whose deviation signal is a scalar anomaly score, not a process-level explanation.

The cross-building generalisation problem is the binding constraint that prevents most existing HVAC FDD work from being deployed in the field. The canonical reference is Zhu and colleagues' study of fault-detection migration between two different building chillers, where conventional supervised classifiers dropped below 55\% target-domain accuracy despite reaching above 95\% on the source, while a domain-adversarial method recovered to 81.27\% on the first operating cycle and 74.93\% on the second \cite{zhu2021chillertransfer}. The closest precedent for a labelled-source-to-unlabelled-target plan of the kind we propose is Ghalamsiah and colleagues' contrastive adaptation network applied to ASHRAE-1312 simulation as a labelled source and real AHU measurements as an unlabelled target \cite{ghalamsiah2025canhvac}. Both approaches treat the transferable artefact as a set of learned neural-network weights, and the structural form of the model itself remains the same across source and target. A discovered process model is a different kind of object: it is an inspectable diagram of expected control flow, and transferring it from one building to another is therefore a structurally different operation that the existing transfer-learning literature has not addressed.

\subsection{Process mining: foundations, adjacent successes, and the empty quadrant in building energy}

Process mining was formalised by van der Aalst as a methodology for discovering process models from event logs, checking observed behaviour for conformance, and enhancing the resulting models with additional perspectives \cite{vanderaalst2016processmining}. The methodology has been applied successfully across cyber-physical and sensor-driven domains adjacent to building energy. Vitale and colleagues applied a full process-mining pipeline to fault diagnosis on the RoAD robotic-arm dataset, reaching an F1 score of up to 98.925\% with a conformance-checking time as low as 0.020 seconds \cite{vitale2025roadcps}. The same group then demonstrated that combining process-mining-derived alignment features with downstream dimensionality reduction substantially outperforms pure conformance checking across four cross-domain datasets, with F1 improvements ranging from 30 to 70 points \cite{vitale2025kbsfeatures}. Husom and colleagues' UDAVA pipeline extended this hybrid pattern to industrial Internet-of-Things sensor validation \cite{husom2026udava}. Methodologically the closest analogue to building HVAC is Moradbeikie and colleagues' work on pasteurisation, where continuous multivariate sensor signals are abstracted into discrete activity sequences via hidden Markov models and conformance is checked against HACCP critical control points \cite{moradbeikie2025pasteurization}. Pasteurisation, like HVAC operation, is a continuous-flow process governed by an explicit regulatory standard; the conceptual mapping from HACCP rules to ASHRAE Guideline 36 control sequences is direct. The technical barrier common to all these applications is event abstraction, the conversion of continuous sensor time-series into discrete activity sequences, which was formalised in the methodology paper by Janssen and colleagues \cite{janssen2021sensoreventlog}.

Despite these adjacent successes, building HVAC remains an empty quadrant in the literature. The novelty claim for the present work rests on the convergence of four independent systematic reviews, each approaching the question from a different angle. Bi and colleagues' survey of 211 AI-based HVAC FDD studies contains no process-mining entries in any leaf of its method taxonomy \cite{bi2024aihvacreview}. Akhramovich and colleagues' review of 47 process-mining studies in Industry 4.0 covers manufacturing, internal logistics, mining, alarm management, supply-chain logistics, and product lifecycle management; building applications do not appear \cite{akhramovich2024pmindustry40}. Brzychczy and colleagues' more recent review of 36 process-mining studies on sensor data documents the same absence at finer granularity: fourteen industrial applications, eleven smart-home applications (all motion-sensor occupant-activity recognition), seven healthcare applications, two smart-office applications, two smart-product applications, and one other, with no building HVAC \cite{brzychczy2025pmsensorreview}. The Lawrence Berkeley National Laboratory team's canonical FDD-methods taxonomy from 2017 organises automated fault-detection methods into quantitative model-based, qualitative model-based, and process-history-based categories, and process mining is not present at any leaf \cite{granderson2017fddtoolsurvey}. Four independent angles converge on the same finding: no published work applies event-log process mining to building HVAC fault detection.

\subsection{Rule-based fault detection in practice, and what process mining adds}

Fault detection is not absent from real buildings; it is simply rule-based. The most widely deployed form is the fault-detection capability embedded in commercial building automation systems. Modern platforms ship libraries of pre-programmed fault rules: the Automated Logic WebCTRL system used at the deployment building of the present study, for example, includes a standard library that matches over one hundred named faults in common HVAC equipment.\footnote{Automated Logic, \emph{Fault Detection and Diagnostics}, \url{https://www.automatedlogic.com/en/products/webctrl-building-automation-system/integration/fdd-reporting-and-dashboards/}.} The characterisation survey of fourteen commercial FDD tools by Granderson and colleagues confirms that this rule-based and threshold-based style is the norm across the commercial landscape \cite{granderson2017fddtoolsurvey}. The strength of a rule library is that it works immediately on connection and requires no baseline data; its structural limitation is that it can only detect the fault patterns that were coded into it in advance, and so it misses novel failure modes and the gradual, multi-variable, sequence-level deviations that no rule author anticipated.

The same rule-based philosophy characterises the most directly comparable research programme, the Lawrence Berkeley National Laboratory team's automated fault-correction line. Pritoni and colleagues field-tested seven hand-coded correction algorithms across 225 distinct equipment items in four buildings on three building-automation-system platforms, targeting common control problems such as incorrectly programmed schedules, unreleased manual overrides, and missed setpoint resets \cite{pritoni2022faultcorrection}. Each algorithm is a pre-encoded rule matched against BAS data and applied through a two-way BACnet writeback to the building's control system. Both the commercial libraries and the LBNL correction algorithms share the assumption that the faults of interest can be enumerated and encoded ahead of time.

ProcessHeal makes the opposite assumption and is therefore complementary to both. Rather than matching observed data against a fixed catalogue of known faults, it learns the expected control flow from the building's own healthy operating data and flags any departure from it, whether or not anyone anticipated that particular departure. The deviation signal is a conformance alignment with formal Petri-net semantics rather than a triggered rule, and it is localised to a specific transition in the discovered model rather than to a named rule. A practical consequence is that ProcessHeal can run on the same trend data that a commercial system such as WebCTRL already records, catching the behaviour-based and sequence-level faults that fall outside the rule library and below the alarm threshold, while the rule library continues to handle the well-understood patterns it was built for. In principle the conformance traces ProcessHeal produces could even serve as input to the LBNL correction algorithms, extending automated correction beyond the patterns pre-encoded by hand. The LBNL FDD dataset family \cite{granderson2020lbnlfdddata} provides the labelled simulation testbed that supports this complementary positioning.

One adjacent body of work must be cited and explicitly distinguished, because its surface vocabulary is similar. Deabes, Bouazza, and Algthami designed a fuzzy Petri-net controller for HVAC temperature regulation, hand-engineering both the Petri-net structure and the fuzzy rules and using the model as a forward control device for an air-conditioning compressor \cite{deabes2023fuzzyptn}. Their work contains no event log, no process-discovery algorithm, no alignment-based conformance checking, and no fault diagnosis. The present work does the methodological opposite: a Petri net is discovered from a real event log of HVAC operation, then used as a reference against which alignment-based conformance checking diagnoses operational deviations. Hand-designed Petri-net controllers and data-discovered Petri-net diagnostic models are two distinct uses of the Petri-net formalism, and they should not be conflated.

\bibliographystyle{plain}
\bibliography{references}

\end{document}
